% Created 2024-03-15 ven. 23:40
% Intended LaTeX compiler: pdflatex
\documentclass[11pt]{article}
\usepackage[utf8]{inputenc}
\usepackage[T1]{fontenc}
\usepackage{graphicx}
\usepackage{longtable}
\usepackage{wrapfig}
\usepackage{rotating}
\usepackage[normalem]{ulem}
\usepackage{amsmath}
\usepackage{amssymb}
\usepackage{capt-of}
\usepackage{hyperref}
\date{}
\title{discours présentation R-GTMG}
\hypersetup{
 pdfauthor={},
 pdftitle={discours présentation R-GTMG},
 pdfkeywords={},
 pdfsubject={},
 pdfcreator={Emacs 29.2 (Org mode 9.7)},
 pdflang={English}}
\usepackage{biblatex}

\begin{document}

\maketitle
Bonjour, aujourd'hui, je vous présente mes travaux sur la réévaluation
de modèle basée sur les règles de transformation de graphe
\section{Introduction}
\label{sec:orgded3f04}

Modéliser des objets complexes, c'est un processus souvent long où on enchaîne
de nombreuses opérations de modélisation.

Aujourd'hui, tous les modeleurs enregistrent ces processus dans ce qui s'appelle
une spécification paramétrique.

Le but étant de les rejouer après avoir modifié, éventuellement, certains
paramètres topologiques ou géométriques, voire ajouté, supprimé ou déplacé
des opérations.

Voyons maintenant un exemple de spécification paramétrique:
\begin{itemize}
	\item Une première opération crée la face carrée f1 et ses arêtes dont e1.
	\item La suivante extrude f1 en un cube et e1 en la face f2.
	\item La triangulation de f2 crée plusieurs faces dont f7.
	\item La dernière opération colore f7.
\end{itemize}

Supposons que l'on modifie cette spécification, on y ajoute l'opération qui
insert un sommet sur l'arête e1 et on la réévalue.

\begin{itemize}
	\item f1 et e1 sont recréées à l'identique
	\item Le sommet inséré scinde e1 en e5 et e6
	\item L'extrusion s'applique sur f1 mais cette fois e5 et e6 sont extrudés en f3 et
	      f4.
\end{itemize}

Et maintenant, comment on réévalue puisque f2 n'existe plus sous sa forme initiale,
pareil pour f7 ?

Intuitivement, f2 est divisé en f3 et f4 alors on voudrait trianguler f3 et/ou f4.

C'est là que l'on a besoin d'un nom persistant pour identifier des cellules par leurs histoire et dire, par exemple, «~f3 est construit comme f2, alors on triangule f3~».

Maintenant que la problématique est posée, on va voir le plan.

\section{plan}
\label{sec:orgdf64a4e}
\begin{itemize}
	\item L'objectif de ces travaux est de permettre la réévaluation de spécifications
	      modifiées notamment avec l'ajout, la suppression et le déplacement
	      d'opérations.
	\item Pour cela, on utilise deux formalismes
	      \begin{itemize}
		      \item Celui des cartes généralisées pour les objets
		      \item Celui des règles Jerboa de transformation de graphe pour les opérations de
		            modélisation
	      \end{itemize}
	\item Cette contribution comprend :
	      \begin{itemize}
		      \item d'une part un mécanisme de nommage persistant basé sur les règles Jerboa
		      \item d'autre part un mécanisme de réévaluation exploitant ce nommage persistant
	      \end{itemize}
\end{itemize}
\section{plan - formalismes}
\label{sec:orgfbb8259}

On commence avec les cartes généralisées, ou G-cartes, pour représenter des
objets
\section{G-cartes}
\label{sec:orgf33c57b}
On prend un objet

Il se décompose en deux volumes connectés par des liaisons 3 en vert

Ses volumes se décomposent en faces connectées par des liaisons 2 en bleu

Qui se décomposent elles-mêmes en arêtes connectées par des liaisons 1 en rouge

Qui se décomposent en sommets connectés par des liaisons 0 en noire

Cette G-carte, c'est un graphe dont les nœuds sont des brins et les arcs des
liaisons.

\section{Orbites}
\label{sec:org99b57a4}

Dans ces G-cartes, on représente les cellules avec des sous-graphes que l'on
appelle orbites.

Une orbite, c'est tous les brins que l'on atteint en suivant certaines liaisons
à partir d'un brin de départ plus les liaisons.

Par exemple~:

Le sommet incident au brin «~a~» contient «~a~», tous les brins atteignables par des
liaisons 1, 2 et 3 et les liaisons.

Cette arête contient le brin «~a~» et tous ceux atteignables par des liaisons 0,2,3.

Une face contient tous les brins atteignables par des liaisons 0,1,3

un volume contient ceux atteignables par des liaisons 0,1,2

On représente aussi toutes sortes d'entités topologiques plus générales que des
cellules comme cette arête de volumes.

\section{Règles de transformation de graphe Jerboa}
\label{sec:org89fe9a5}

On a formalisé les objets avec les G-cartes, maintenant, on formalise les
opérations avec les règles de transformations de graphe Jerboa.

On prend l'exemple de la triangulation de face.

De manière générale, une règle se compose d'une partie gauche qui filtre un
sous-graphe avant la transformation d'un objet et d'une partie droite qui décrit
la transformation.

À gauche, n0 porte le type d'orbite face <0,1,3>. Quand on applique la règle sur
a0, n0 filtre cette face.

À droite, n0 est préservé donc ses brins sont préservés. Les liaisons 1 sont
supprimées, donc on a ce symbole tiret dans n0.

n2 est un nœud créé donc ses brins sont créés par copie de ceux de n0. %
Les liaisons 0 sont renommées en liaisons 1 et les liaisons 1 en liaisons 2. %
Donc, dans l'objet, il y a des liaisons 1 et 2 entre les brins.

Enfin, n1 est créé, on a renommage et suppression.

Une liaison 1 connecte n0 à n1, pareil pour leurs brins.
Même principe entre n1 et n2.

Les liaisons dans les nœuds sont des arcs implicites.
Celles entre les nœuds sont des arcs explicites.

\subsection{Orbites}
\label{sec:orgcf72bd7}

On peut aussi suivre des orbites dans les règles.

Par exemple la face de volume incidente à n0.

À gauche, elle contient n0 avec ses arcs 0,1

À droite, elle contient n0 avec ses arcs 0, l'arc 1 qui le connecte à
n1, l'arc 0 qui connecte n1 à n2 et les arcs 1 de n2.

On remarque qu'il y a une seule orbite face dans la règle, alors que l'on en a 4
dans l'objet. %
C'est dû aux renommages de cet arc implicite 1, à gauche, qui ne sont ni 0
ni 1 à droite, donc on a une scission de la face de volume filtrée par n0.

\subsection{Origines}
\label{sec:org341560c}

En plus, les arcs implicites 0 et 1 qui constituent notre face de volume sont
des renommages de cet arc implicite 0 à gauche, donc cette face de volume a pour
origine cette arête de face de volume de type <0>. %

Dans l'objet, chaque face a pour origine une arête préservée.

Par exemple cette face verte vient de cette arête verte, la bleue de cette bleue et
ainsi de suite.

Grâce à une analyse statique en pré-calcul des règles, on est capable de
détecter tout type d'évènement sur les orbites indépendamment des objets.

Notamment, la scission, fusion, etc.
\section{Plan Nommage Persistant}
\label{sec:orgaa5036e}

Avec ces formalismes, on peut maintenant voir le nommage persistant.

On reprend la spécification de l'introduction.
\section{Nommage Persistant}
\label{sec:org84648da}

\subsection{Nom persistant brin}
\label{sec:org53f1b92}

Ici, on propose d'enregistrer l'histoire des brins comme leurs noms
persistants. %
Chaque brin contient alors les numéros d'opérations et de nœuds qui les créent
ou réécrivent.

Cette première règle crée une face carrée.
Elle a 8 nœuds dont n3 et n6.

n3, crée le brin 3. Son histoire, c'est application 1 nœud n3, donc [1n3].

n6 crée le brin 6, son histoire, c'est [1n6].

Quand on applique l'extrusion sur le carré, on sélectionne le brin 6. %
Sauf que le brin pourrait changer, alors dans la spécification, on met son histoire [1n6].

n1 filtre le brin 3 et n2 et n4 créent ses copies 33 et 35. %
On construit leurs histoires à partir de celle de 3, suivie des leurs.

L'histoire de 33 est [1n3;2n2]. %
L'histoire de 35 est [1n3;2n4].

Ensuite, on applique la triangulation sur la face désignée par le brin 33 et on
enregistre son histoire.

À droite, n0 réécrit 35, comme c'est la troisième opération, son histoire devient
[1n3;2n4;3n0].

Enfin, on colorie la face désignée par 35.

On vient de nommer des brins de manière unique en représentant leurs histoires, il s'agit du nommage persistant des brins.

Mais, dans une spécification, les paramètres topologiques sont des orbites.
Donc, on doit représenter les histoires de ces orbites.

\subsection{Nom persistant orbite}
\label{sec:orgf5219a2}

Pour reconstituer l'histoire d'une orbite, on prend le nom persistant qui la
désigne et on remonte l'histoire décrite.

Cette histoire prend la forme d'un arbre.

On prend l'orbite désignée par [1n3;2n4;3n0] dans l'opération de coloration.

Les arcs implicites de ce n0 désignent une face.

L'histoire commence alors avec une face incidente à n0 dans la troisième opération.

Cette flèche noire dit qu'elle provient de la scission de cette face. %
Cette flèche rouge dit qu'elle a pour origine une arête.

On remonte l'histoire d'un cran à la deuxième opération et son nœud n4.

La face que l'on suit est créé lors de l'extrusion, elle a pour origine une
arête.

Cette arête, elle est créé à partir de l'extrusion d'un brin.

Enfin, on remonte encore l'histoire, cette fois à la toute première opération et son nœud n3.

Cette règle crée le brin et l'arête incidents au nœud 3.

On vient de reconstituer l'histoire d'une orbite statiquement à partir d'un nom
persistant et des règles impliquées.

On suit la même procédure pour chacun des paramètres topologiques enregistrés dans une spécification et uniquement ceux-ci.

\section{Plan réévaluation}
\label{sec:org04c39be}

Grâce aux histoires que l'on reconstitue, on peut réévaluer une spécification.

\section{Réévaluation}
\label{sec:org3b38850}

On reprend l'histoire précédente, cette fois de haut en bas, et on va réévaluer
la spécification modifiée. %

La première opération recrée la face carrée à l'identique.

n3 crée un unique brin, 3. On retrouve nos orbites arête et brin.

L'opération suivante a été ajoutée, insert un sommet sur l'arête désignée par
le brin 3.

Résultat, cette modification de spécification scinde l'arête en deux, l'une
désignée par le brin 3, l'autre désignée par le brin 4.

On a donc une histoire par arête, les brins eux, restent inchangés.

Pour l'opération d'extrusion, son paramètre topologique a été déterminé au
préalable dans un autre arbre, exactement comme on est en train de faire.

Maintenant, on a deux histoires que l'on peut suivre~: une pour chaque arête.

Chaque arête est extrudée en une nouvelle face, et chaque des brins en une
arête.

Ici, n4 crée plusieurs brins mais une seule copie de 3, qui est 37, et une seule
copie de 4, qui est 52.

Enfin, on peut réappliquer la triangulation sur chaque branche.

n0 réécrit 37 et 52 à l'identique.

C'est la dernière application enregistrée, 37 et 52 désignent les faces à colorer.

\section{Conclusion}
\label{sec:orga7a7cfd}

Pour conclure, on propose un mécanisme de réévaluation adapté aux systèmes de
modélisation à base de règles.

Pour cela, on met place un nommage persistant différents niveaux~:

un pour les brins, ce qui rend le nommage unique et homogène en toutes
dimensions~;

l'autre pour les orbites avec une histoire complète au cours de laquelle on peut
suivre l'évolution d'une orbite et de ses origines

Le nommage d'une orbite est fait entièrement à partir des informations obtenues lors de
l'analyse statique des règles en pré-calcul.

Pendant la réévaluation, on suit toutes les histoires possibles ce qui permet de
mettre en place des stratégies comme la décision de réappliquer des opérations
sur une ou plusieurs branches.

Actuellement, on travaille sur la réévaluation de scripts qui appellent des
règles dans les structures de contrôles habituelles comme les alternatives et
les boucles.
\end{document}
